\section{Measuring the dependence on wording}\label{sec:task}

\subsection{What is held fixed}

We evaluate on COCO val2017 \citep{lin2014microsoft}, 5{,}000 images of which \NIMAGES{}
carry annotations, over 80 categories, and on the LVIS v1.0 minival split
\citep{gupta2019lvis}, 4{,}809 images and 1{,}203 categories; LVIS images are the same COCO 2017 JPEGs, so the two corpora share a pixel
source and differ in annotations and vocabulary.

The detector is Grounding DINO \citep{liu2024grounding} with the public Swin-T checkpoint.
It is used exactly as distributed: no fine-tuning, no prompt tuning, no threshold tuning.
For each image we query the categories that are actually annotated in it, which is the
realistic case for an annotation assistant (a user points at what is there) and avoids
inflating the false-positive count with categories that were never asked for.

Each image is queried five times. In the $v$-th run the phrase for category $c$ is

\begin{equation}
\phi_v(c) = \big(\text{name},\ \text{a name},\ \text{a photo of a name},\
\text{an image of a name},\ \text{names}\big)_v ,
\end{equation}

that is: the bare name, an indefinite article, the two templates that circulate as prompt
engineering for region--text models, and the plural. We call these the five \emph{wordings}
and write $\phi$ for a phrase. Everything else, including the image, the category, the
post-processing and the weights, is identical across the five runs, so any difference in
the output is attributable to the phrase.

\subsection{Detection records}

Within a run, detections are kept at a score threshold of $0.01$ and capped at the top 80
per (image, category) pair, which is well inside the operating range of the detector and
leaves \NDETS{} records on COCO. A detection is \emph{correct} if it matches a
ground-truth instance of the queried category at IoU $\geq 0.5$, assigned greedily in
descending score. Matching is done separately inside each run: the five runs are five
queries, and letting them share one ground-truth instance would score four of the five as
wrong by construction.

Table~\ref{tab:setup} summarises the two corpora. The band $p\in[0.3,0.9)$ is where a human
reviews boxes: it is \BANDFRAC{} of the output of a run and holds \BANDPOS{} of its true
positives, and it is where the detector is least reliable, with a band error of
\BANDIDN.

\begin{table}[t]
\caption{The five wordings on the same images with the same weights. A row is one phrasing; nothing but the text query changes between rows. ``Band'' is $p\in[0.3,0.9)$ of the uncalibrated score.}
\label{tab:setup}
\begin{tabular*}{\textwidth}{@{\extracolsep{\fill}}lrrrrrr@{}}
\toprule
Wording & detections & per image & positive rate & AP & share in band & ECE \\
\midrule
bare               & 1{,}063{,}949 & 215 & 3.2\% & 0.7770 & 4.3\% & 0.0500 \\
article            & 1{,}155{,}657 & 233 & 2.8\% & 0.7430 & 2.0\% & 0.0602 \\
photo              & 1{,}170{,}394 & 236 & 2.6\% & 0.4622 & 1.9\% & 0.0596 \\
image              & 1{,}169{,}850 & 236 & 2.7\% & 0.5020 & 2.2\% & 0.0512 \\
plural             & 1{,}133{,}371 & 229 & 3.0\% & 0.7217 & 2.6\% & 0.0474 \\
\bottomrule
\end{tabular*}
\end{table}


\subsection{Metrics, and why they are computed per run}

The five runs are evaluated separately and the results averaged; pooling them would let one
phrasing contribute five copies of the same ground-truth instance and would hide exactly the
effect we are measuring. Inside a run we report:

\begin{itemize}
\item the expected calibration error over 15 bins (ECE), and the same quantity restricted to
the operating band with 10 bins (band ECE);
\item the Brier score, a proper score, and its restriction to the band (band Brier). The
restricted Brier is the quantity we trust most: it needs no binning, and on the band it is
computed on the boxes a reviewer would actually see;
\item average precision (AP) at 101 recall points, computed per category, and the area under
the risk--coverage curve (AURC) for the selective-prediction view.
\end{itemize}

AP here is not the official benchmark metric and is not comparable to published
zero-shot numbers: matching happens inside a run and only over the categories that were
asked for, so there is no penalty for the categories the image does not contain. It is
used only to compare the five runs, and the calibrators within a run, on identical
candidates.

The band mask of a run is fixed by that run's \emph{uncalibrated} score, so every calibrator
is scored on the same detections. Calibrators are fitted on one half of the images and
evaluated on the other, split by image rather than by detection: one image contributes about
\DETPERIMG{} queries that are strongly correlated, and a detection-level split would place
near-copies on both sides.

\subsection{The quantity we are after}

For an object, we take the highest-scoring box the detector returns for its category in each
of the five runs. If the phrasing did not matter, these five numbers would agree. We report
their range, the fraction of objects whose range exceeds $0.1$, and how often the five runs
disagree about whether the box is correct at all. The same object, five scores, is the unit
of the measurement; Section~\ref{sec:exp-measure} reports it.
